\chapter{2020年天津卷}

\section{单选题}

\begin{problem}
设全集 \(U=\{-3,-2,-1,0,1,2,3\}\)，集合 \(A=\{-1,0,1,2\}\)，\(B=\{-3,0,2,3\}\)，则 \(A\cap\left( U\setminus B \right)=\)（\quad）

\choices
    {\(\{-3,3\}\)}
    {\(\{0,2\}\)}
    {\(\{-1,1\}\)}
    {\(\{-3,-2,-1,1,3\}\)}
\end{problem}

\begin{problem}
设 \(a\in\R\)，则“\(a>1\)”是“\(a^2>a\)”的（\quad）

\choices
    {充分不必要条件}
    {必要不充分条件}
    {充要条件}
    {既不充分也不必要条件}
\end{problem}

\begin{problem}
函数 \(y=\dfrac{4x}{x^2+1}\) 的图象大致为（\quad）

\choices
    {\choicebitmap[width=0.15\paperwidth]{img/2020/tianjin/q3_fig1.png}}
    {\choicebitmap[width=0.15\paperwidth]{img/2020/tianjin/q3_fig2.png}}
    {\choicebitmap[width=0.15\paperwidth]{img/2020/tianjin/q3_fig3.png}}
    {\choicebitmap[width=0.15\paperwidth]{img/2020/tianjin/q3_fig4.png}}
\end{problem}

\begin{problem}
从一批零件中抽取 80 个，测量其直径（单位：mm），将所得数据分为 9 组 \([5.31,5.33)\)，\([5.33,5.35)\)，\(\cdots\)，\([5.45,5.47)\)，\([5.47,5.49]\)，并整理得到如下频率分布直方图，则在被抽取的零件中，直径落在区间 \([5.43,5.47)\) 内的个数为（\quad）

\begin{center}
\bitmapfigure[max width=0.90\linewidth]{img/2020/tianjin/q4_fig1.png}
\end{center}

\choices
    {10}
    {18}
    {20}
    {36}
\end{problem}

\begin{problem}
若棱长为 \(2\sqrt3\) 的正方体的顶点都在同一球面上，则该球的表面积为（\quad）

\choices
    {\(12\pi\)}
    {\(24\pi\)}
    {\(36\pi\)}
    {\(144\pi\)}
\end{problem}

\begin{problem}
设 \(a=3^{0.7}\)，\(b=\left( \dfrac13 \right)^{-0.8}\)，\(c=\log_{0.7}0.8\)，则 \(a\)，\(b\)，\(c\) 的大小关系为（\quad）

\choices
    {\(a<b<c\)}
    {\(b<a<c\)}
    {\(b<c<a\)}
    {\(c<a<b\)}
\end{problem}

\begin{problem}
设双曲线 \(C\) 的方程为 \(\dfrac{x^2}{a^2}-\dfrac{y^2}{b^2}=1\)（\(a>0\)，\(b>0\)），过抛物线 \(y^2=4x\) 的焦点和点 \((0,b)\) 的直线为 \(l\). 若 \(C\) 的一条渐近线与 \(l\) 平行，另一条渐近线与 \(l\) 垂直，则双曲线 \(C\) 的方程为（\quad）

\choices
    {\(\dfrac{x^2}{4}-\dfrac{y^2}{4}=1\)}
    {\(x^2-\dfrac{y^2}{4}=1\)}
    {\(\dfrac{x^2}{4}-y^2=1\)}
    {\(x^2-y^2=1\)}
\end{problem}

\begin{problem}
已知函数 \(f(x)=\sin\left( x+\dfrac{\pi}{3} \right)\). 给出下列结论：

\begin{circlelist}
\item \(f(x)\) 的最小正周期为 \(2\pi\)；

\item \(f\left( \dfrac{\pi}{2} \right)\) 是 \(f(x)\) 的最大值；

\item 把函数 \(y=\sin x\) 的图象上所有点向左平移 \(\dfrac{\pi}{3}\) 个单位长度，可得到函数 \(y=f(x)\) 的图象.
\end{circlelist}

其中所有正确结论的序号是（\quad）

\choices
    {\(\circled{1}\)}
    {\(\circled{1}\circled{3}\)}
    {\(\circled{2}\circled{3}\)}
    {\(\circled{1}\circled{2}\circled{3}\)}
\end{problem}

\begin{problem}
已知函数 \(f(x)=x^3\)（\(x\ge 0\)），\(f(x)=-x\)（\(x<0\)）. 若函数 \(g(x)=f(x)-\left| kx^2-2x \right|\)（\(k\in\R\)）恰有 4 个零点，则 \(k\) 的取值范围是（\quad）

\choices
    {\(\left( -\infty,-\dfrac12 \right)\cup\left( 2\sqrt2,+\infty \right)\)}
    {\(\left( -\infty,-\dfrac12 \right)\cup\left( 0,2\sqrt2 \right)\)}
    {\(\left( -\infty,0 \right)\cup\left( 0,2\sqrt2 \right)\)}
    {\(\left( -\infty,0 \right)\cup\left( 2\sqrt2,+\infty \right)\)}
\end{problem}

\section{填空题}

\begin{problem}
\(\i\) 是虚数单位，复数 \(\dfrac{8-\i}{2+\i}=\fillinblank{}\).
\end{problem}

\begin{problem}
在 \(\left( x+\dfrac{2}{x^2} \right)^5\) 的展开式中，\(x^2\) 项的系数是\fillinblank{}.
\end{problem}

\begin{problem}
已知直线 \(x-\sqrt3 y+8=0\) 和圆 \(x^2+y^2=r^2\)（\(r>0\)）相交于 \(A\)，\(B\) 两点. 若 \(|AB|=6\)，则 \(r=\fillinblank{}\).
\end{problem}

\begin{problem}
已知甲，乙两球落入盒子的概率分别为 \(\dfrac12\) 和 \(\dfrac13\). 假定两球是否落入盒子互不影响，则甲，乙两球都落入盒子的概率为\fillinblank{}；甲，乙两球至少有一个落入盒子的概率为\fillinblank{}.
\end{problem}

\begin{problem}
已知 \(a>0\)，\(b>0\)，且 \(ab=1\)，则 \(\dfrac1{2a}+\dfrac1{2b}+\dfrac8{a+b}\) 的最小值为\fillinblank{}.
\end{problem}

\begin{problem}
如图，在四边形 \(ABCD\) 中，\(\angle B=60^\circ\)，\(|AB|=3\)，\(|BC|=6\)，且 \(\overrightarrow{AD}=\lambda\overrightarrow{BC}\)，\(\overrightarrow{AD}\cdot\overrightarrow{AB}=-\dfrac32\)，则实数 \(\lambda\) 的值为\fillinblank{}，若 \(M\)，\(N\) 是线段 \(BC\) 上的动点，且 \(|MN|=1\)，则 \(\overrightarrow{DM}\cdot\overrightarrow{DN}\) 的最小值为\fillinblank{}.

\begin{center}
\bitmapfigure[max width=0.76\linewidth]{img/2020/tianjin/q15_fig1.png}
\end{center}
\end{problem}

\section{解答题}

\begin{problem}
在 \(\triangle ABC\) 中，角 \(A\)，\(B\)，\(C\) 所对的边分别为 \(a\)，\(b\)，\(c\). 已知 \(a=2\sqrt2\)，\(b=5\)，\(c=\sqrt{13}\).

\begin{enumerate}
\item 求角 \(C\) 的大小；
\item 求 \(\sin A\) 的值；
\item 求 \(\sin\left( 2A+\dfrac{\pi}{4} \right)\) 的值.
\end{enumerate}
\end{problem}

\begin{problem}
如图，在三棱柱 \(ABC-A_1B_1C_1\) 中，\(CC_1\perp\) 平面 \(ABC\)，\(AC\perp BC\)，\(AC=BC=2\)，\(CC_1=3\)，点 \(D\)，\(E\) 分别在棱 \(AA_1\) 和棱 \(CC_1\) 上，且 \(AD=1\)，\(CE=2\)，\(M\) 为棱 \(A_1B_1\) 的中点.



\begin{enumerate}
\item 求证：\(C_1M\perp B_1D\)；
\item 求二面角 \(B-B_1E-D\) 的正弦值；
\item 求直线 \(AB\) 与平面 \(DB_1E\) 所成角的正弦值.
\end{enumerate}

\bitmapfigure[max width=0.80\linewidth]{img/2020/tianjin/q17_fig1.png}
\end{problem}

\begin{problem}
已知椭圆 \(\dfrac{x^2}{a^2}+\dfrac{y^2}{b^2}=1\)（\(a>b>0\)）的一个顶点为 \(A(0,-3)\)，右焦点为 \(F\)，且 \(|OA|=|OF|\)，其中 \(O\) 为原点.

\begin{enumerate}
\item 求椭圆方程；
\item 已知点 \(C\) 满足 \(3\overrightarrow{OC}=\overrightarrow{OF}\)，点 \(B\) 在椭圆上（\(B\) 异于椭圆的顶点），直线 \(AB\) 与以 \(C\) 为圆心的圆相切于点 \(P\)，且 \(P\) 为线段 \(AB\) 的中点. 求直线 \(AB\) 的方程.
\end{enumerate}
\end{problem}

\begin{problem}
已知 \(\{a_n\}\) 为等差数列，\(\{b_n\}\) 为等比数列，\(a_1=b_1=1\)，\(a_5=5\left( a_4-a_3 \right)\)，\(b_5=4\left( b_4-b_3 \right)\).

\begin{enumerate}
\item 求 \(\{a_n\}\) 和 \(\{b_n\}\) 的通项公式；
\item 记 \(\{a_n\}\) 的前 \(n\) 项和为 \(S_n\)，求证：\(S_nS_{n+2}<S_{n+1}^2\)（\(n\in\N^*\)）；
\item 对任意的正整数 \(n\)，设 \(c_n=\dfrac{(3a_n-2)b_n}{a_na_{n+2}}\)（\(n\) 为奇数），\(c_n=\dfrac{a_{n-1}}{b_{n+1}}\)（\(n\) 为偶数），求数列 \(\{c_n\}\) 的前 \(2n\) 项和.
\end{enumerate}
\end{problem}

\begin{problem}
已知函数 \(f(x)=x^3+k\ln x\)（\(k\in\R\)），\(f'(x)\) 为 \(f(x)\) 的导函数.

\begin{enumerate}
\item 当 \(k=6\) 时，
\begin{enumerate}
\item 求曲线 \(y=f(x)\) 在点 \((1,f(1))\) 处的切线方程；
\item 求函数 \(g(x)=f(x)-f'(x)+\dfrac9x\) 的单调区间和极值；
\end{enumerate}
\item 当 \(k\ge-3\) 时，求证：对任意的 \(x_1\)，\(x_2\in[1,+\infty)\)，且 \(x_1>x_2\)，有 \(\dfrac{f'(x_1)+f'(x_2)}{2}>\dfrac{f(x_1)-f(x_2)}{x_1-x_2}\).
\end{enumerate}
\end{problem}

